%%% BAC exercise template %%%
% template modified by Derek Huang, original by Sean Cox.
\documentclass{article}
\usepackage[margin=1in]{geometry}
\usepackage{amsmath, amssymb, amsfonts, enumitem, fancyhdr, tikz}
% get rid of paragraph indent
\setlength{\parindent}{0 pt}
% allow section.equation numbering
\numberwithin{equation}{section}
% allows you to copy-paste code segments. requires pygments, which can be
% installed from PyPI with pip install Pygments.
% warning: minted does NOT work with Texmaker if you have TeX Live installed
% on WSL but Texmaker installed natively on Windows!
%\usepackage{minted}
% alternative to minted that does not require Python, LaTeX only. listings is
% however disgusting out of the box and some setup is required.
\usepackage{listings, xcolor}
% makes clickable links to sections
\usepackage{hyperref}
% make the link colors blue, as well as cite colors. urls are magenta
\hypersetup{
    colorlinks, linkcolor = blue, citecolor = blue, urlcolor = magenta
}
% fancy pagestyle so we can use fancyhdr for fancy headers/footers
\pagestyle{fancy}
% add logo in right of header. note that you will have to adjust logo path!
\fancyhead[R]{\includegraphics[scale = 0.15]{../bac_logo1.png}}
% don't show anything in the left and center header
\fancyhead[L, C]{}
% give enough space for logo by reducing top margin height, head separator,
% increasing headerheight. see Figure 1 in the fancyhdr documentation. if
% \topmargin + \headheight + \headsep = 0, original text margins unchanged.
\setlength{\topmargin}{-60 pt}
\setlength{\headheight}{50 pt}
\setlength{\headsep}{10 pt}
% remove decorative line in the fancy header
\renewcommand{\headrulewidth}{0 pt}

% color definitions for listings syntax highlighting. uses colors borrowed
% from the VS Code Dark+ and Abyss standard themes.
\definecolor{KwColor}{RGB}{153, 102, 184}     % keyword color
\definecolor{VarColor}{RGB}{86, 156, 214}     % variables/identifier color
\definecolor{StrColor}{RGB}{209, 105, 105}    % string color
\definecolor{CmtColor}{RGB}{106, 153, 85}     % comment color

% general listings configuration for all languages
\lstset{
    % change keyword, identifier, comment, string colors
    keywordstyle = \color{KwColor},
    commentstyle = \color{CmtColor},
    identifierstyle = \color{VarColor},
    stringstyle = \color{StrColor},
    % no spaces in strings
    showstringspaces = false,
    % monospace font by default
    basicstyle = \ttfamily,
    % tabsize 8 by default, this is not the 1960s
    tabsize = 4,
    % add line numbers to the left with gray typewriter font
    numbers = left,
    numberstyle = \color{gray}\ttfamily,
    % change distance from code block from 10 pt to 5 pt
    numbersep = 5 pt
}

% title, author + thanks, date
\title{Exercise 5}
\author{Derek Huang\thanks{NYU Stern 2021, BAC Advanced Team.}}
\date{March 30, 2021}

% shortcut links. the % characters strip extra spacing.
\newcommand{\pytest}{\href{https://docs.pytest.org/en/stable/}{pytest}}
\newcommand{\minimize}{%
    \href{%
        https://docs.scipy.org/doc/scipy/reference/generated/%
        scipy.optimize.minimize.html%
    }{\texttt{scipy.optimize.minimize}}%
}
% don't wanna keep typing this out
\newcommand{\tr}{\operatorname{tr}}
\newcommand{\npinv}{%
    \href{%
        https://numpy.org/doc/stable/reference/generated/numpy.linalg.inv.html%
    }{\texttt{numpy.linalg.inv}}%
}
\newcommand{\pdb}{%
    \href{https://docs.python.org/3/library/pdb.html}{\texttt{pdb}}%
}

\begin{document}

\maketitle
% need to include this after making title to undo the automatic
% \thispagestyle{plain} command that is issued.
\thispagestyle{fancy}


\section{Introduction}

The goal of this exercise is to implement a linear discriminant classifier,
computing the sample covariance matrix using a maximum likelihood estimate
with optional shrinkage.

\section{Instructions}

\subsection{General}

The \texttt{exercise\_05.py} file contains a skeleton for the
\texttt{LinearDiscriminantAnalysis} class, unit tests, and \pytest{} fixtures.
Your job is to finish implementing the \texttt{LinearDiscriminantAnalysis}
class's \texttt{fit} and \texttt{score} methods by replacing the
\texttt{\#\#\# your code goes here \#\#\#} comment blocks with appropriate
Python code.

\medskip

Your code \textbf{must} be written in the areas marked by these blocks. Do
\textbf{not} change any of the pre-written code. The exercise is complete
$ \Leftrightarrow $ \texttt{pytest /path/to/exercise\_05.py} executes with
zero test failures. Note that attributes ending with \texttt{\_}, ex.
\texttt{coef\_}, are created \textbf{during} the fitting process.

\subsection{Implementing \texttt{fit}}

The class prior estimates must be assigned to the local variable
\texttt{priors}, the class mean estimates must be assigned to the local
variable \texttt{means}, and the covariance matrix estimate must be assigned
to the local variable \texttt{cov}. At the end of the \texttt{fit} method,
\texttt{priors} is assigned to \texttt{self.priors\_}, \texttt{means} is
assigned to \texttt{self.means\_}, and \texttt{cov} is assigned to
\texttt{self.covariance\_}.

\subsection{%
    Computing \texttt{priors\_}, \texttt{means\_}, \texttt{covariance\_}%
}

The estimates for the class priors and class means should be computed via
maximum likelihood. Note that the shape of \texttt{means\_} is specified in
the docstring to be \texttt{(n\_classes, n\_features)}. The covariance matrix
estimate $ \hat{\mathbf{\Sigma}} \succeq \mathbf{0} \in
\mathbb{R}^{d \times d} $ is such that for $ \alpha \in [0, 1] $ we have
\begin{equation} \label{cov_est}
    \hat{\mathbf{\Sigma}} \triangleq \alpha\frac{\tr(\mathbf{S})}{d}\mathbf{I}
        + (1 - \alpha)\mathbf{S}
\end{equation}

Here $ \mathbf{S} \succeq \mathbf{0} \in \mathbb{R}^{d \times d} $ is the
maximum likelihood estimate for the covariance matrix and
$ \tr(\mathbf{S}) $ is its trace, i.e.
\begin{equation*}
    \tr(\mathbf{S}) = \sum_{i = 1}^ds_{ii}
\end{equation*}

Note that in this context, $ \tr(\mathbf{S}) / d $ is the average of the
input features' variances. In practice, other \textit{shrinkage estimators}
may be used, especially when $ d $ is much larger than $ N $, for example
the \href{http://www.ledoit.net/ole1a.pdf}{Ledoit-Wolf}\footnote{
    In portfolio optimization, Ledoit-Wolf instead refers to
    \href{http://www.ledoit.net/honey.pdf}{%
        a particular constant correlation shrinkage estimator%
    }.
} and
\href{https://arxiv.org/pdf/0907.4698.pdf}{%
    oracle approximating shrinkage%
} estimators, which are beyond the scope of our introductory coverage.

\subsection{Implementing \texttt{score}}

The \texttt{score} method of the \texttt{LinearDiscriminantAnalysis} class
computes the accuracy of the predictions. Given parameters $ \theta $,
classification function $ F_\theta $ (the classifer), and data
$ \mathcal{D} \triangleq \{(\mathbf{x}_1, y_1), \ldots (\mathbf{x}_N, y_N)\} $,
where the $ y_k $ values may either be numeric or class labels, the accuracy
$ \mathcal{A}_\theta $ of the model on the data $ \mathcal{D} $ is such that
\begin{equation*}
    \mathcal{A}_\theta(\mathcal{D}) = \frac{1}{|\mathcal{D}|}
        \sum_{(\mathbf{x}, y) \in \mathcal{D}}\mathbb{I}_{\{y\}}\left(
            F_\theta(\mathbf{x})
        \right)
\end{equation*}

Note that $ \operatorname{im}\mathcal{A}_\theta = [0, 1] $, so for
data $ \mathcal{D} $, $ \mathcal{A}_\theta(\mathcal{D}) $ is the
fraction of examples the model correctly clasified.

\subsection{Tips}

\begin{enumerate}
    \item
    Carefully read the \texttt{LinearDiscriminantAnalysis} docstrings and
    review slide 8.

    \item
    In the \texttt{fit} method, don't forget to either assign your covariance
    matrix estimate to a local variable named \texttt{cov} or replace the
    variable name passed to the \npinv{} call with the name of the local
    variable holding a reference to the covariance matrix estimate.

    \item
    In the \texttt{score} method, you should first call the \texttt{predict}
    method and use it get predicted labels.    

    \item
    Use NumPy functions whenever possible for efficiency, brevity, and ease
    of debugging.

    \item
    Invoke \pytest{} with the \texttt{--pdb} flag to start the \pdb, the
    Python debugger. Doing so allows you to inspect the variables in the
    current call frame and look more closely at what went wrong.
\end{enumerate}




\end{document}